\documentclass[11pt]{article}

\usepackage[margin=1.1in]{geometry}
\usepackage{amsmath}
\usepackage{amssymb}
\usepackage{booktabs}
\usepackage{microtype}
\usepackage{xcolor}
\usepackage{caption}
\usepackage{array}
\usepackage[colorlinks=true, linkcolor=blue!50!black, citecolor=blue!50!black, urlcolor=blue!50!black]{hyperref}

\captionsetup{font=small, labelfont=bf}

\newcommand{\code}[1]{\texttt{#1}}

\title{\textbf{The Inert Axis: When a Style Parameter Is\\ Wired, Swept, and Never Reached}}
\author{Allen Hsieh\\[2pt] {\small FishAI project}}
\date{August 2026}

\begin{document}

\maketitle

\begin{abstract}
\noindent
The FishAI bot laboratory ships a roster of nine play styles for the \code{us54} variant of the
card game Literature (Canadian Fish), named along an aggressive-to-conservative axis whose defining knob is
\code{declareThreshold}, the confidence a policy demands before a speculative declare. Direct
decision-level instrumentation --- replaying real games and handing two policy vectors the
identical observation and random seed at every decision point --- shows that across the range the
roster actually spans --- every shipped bar sits at $0.775$ or above, the span being
$[0.775, 0.98]$ --- this knob changes \emph{nothing}: zero divergent decisions at every measured
threshold from $0.775$ up (through $0.975$ at roster stage, and through $1$ in the pilot), in
populations of up to 275{,}380 decisions. The mechanism is not an unreachable code
path. The speculative-declare path fires constantly (171 speculative to 138 certain declares in
one 40-game population); rather, the shared inference engine's confidence estimates are
\emph{bimodal} --- a plan sits either well above ${\sim}0.975$ or well below ${\sim}0.775$, almost
never between --- so every threshold inside the band selects the identical set of declares. We
describe the instrument, the finding, the mechanism, and the repairs: a relocated behavioural
gate, a relabelled narrative, a regression test asserting that no roster vector is inert, and a
search-side statistic that counts byte-identical candidates. The general lesson is
methodological: parameter sweeps that report only outcomes cannot distinguish ``no effect'' from
``never consulted.'' Decision-level divergence is the honest unit.
\end{abstract}

\section{Introduction}
\label{sec:intro}

A common failure mode in studies of parameterised game-playing agents is quiet: a parameter is
declared, plumbed through the code, swept across a grid, and reported in a results table --- and at
no point in any of those steps does the software ever act differently because of it. The sweep
returns flat outcomes; the flatness is read as ``this axis does not matter much''; and the
conclusion, while accidentally true in its operational content, badly misstates \emph{why}. The
knob was not weak. It was dead.

This paper documents a fully worked instance of that failure, found and repaired inside the
FishAI project --- an engine, bot laboratory, and results site for the six-player partnership card
game Literature (Canadian Fish; colloquially just ``Fish''), in its 54-card ``\code{us54}'' variant
\cite{rules}. The laboratory's nine-style roster \cite{styles, v05} was designed and \emph{named}
around a single axis: how much confidence a style demands before making a speculative declare,
the knob \code{declareThreshold}. Aggressive styles were given low bars, conservative styles high
ones, and the roster's public narrative was a spectrum from one to the other.

Measured at the level of individual decisions, that axis turns out to be inert across the entire
range the roster spans. Every measured value of \code{declareThreshold} from $0.775$ through
$0.975$ --- and through $1$ in the pilot, covering the shipped $[0.775, 0.98]$ span ---
produces \emph{byte-identical play}: not statistically indistinguishable play, but the same
action, at every decision point, in every game, on every seed. The styles do remain genuinely
distinct --- between $0.39\%$ and $2.89\%$ of their decisions differ from the control's --- but
through other knobs entirely, not along the axis their labels advertise.

The finding has three parts, and the paper is organised around them:

\begin{enumerate}
  \item \textbf{The instrument} (Section~\ref{sec:instrument}): a decision-level divergence probe
  that replays real games and, at every decision point, hands two policy vectors the identical
  observation (\code{SeatView}) and identical random seed, then compares their two actions
  exactly. Because the policy layer is a pure function of (view, policy, seed), a count of zero
  divergent decisions is a certificate of behavioural identity, not a failed significance test.
  \item \textbf{The finding and its mechanism} (Sections~\ref{sec:finding}
  and~\ref{sec:mechanism}): the zero-divergence result, its attribution one knob at a time, the
  threshold sweeps, and the reason --- the inference engine's confidence estimates for
  speculative declare plans are bimodal, because certainty usually preempts speculation.
  \item \textbf{The consequences} (Section~\ref{sec:consequences}): a behavioural gate relocated
  so the one style built on it exists at all; the project's own labelling corrected in place
  rather than quietly rewritten; a regression test that permanently asserts no roster vector is
  inert; and a search-side statistic, \code{inertCandidates}, that measures the same phenomenon
  from the other direction.
\end{enumerate}

Section~\ref{sec:lessons} draws the general lessons for parameter studies of agents, and
Section~\ref{sec:limitations} states what this result does not show.

\section{Setting: the game, the rule change, and the roster}
\label{sec:setting}

\subsection{The \code{us54} declare rule}

Fish is a six-player, two-team game of asking opponents for specific cards and \emph{declaring}
complete six-card sets (``books'') by naming the exact holder of every card. FishAI plays two
rule sets. Under the 48-card default (\code{pagat48}) the error outcomes split: declaring a set
an opponent holds a card of awards it to the opponents, while misassigning a set one's own team
wholly holds \emph{voids} it --- nobody scores. The \code{us54} variant collapses the split in one
decisive row of its decision table --- row~14: a declare with \emph{any error at all} awards the
set to the \emph{opposing} team; the void outcome is abolished \cite{rules}. Combined with
\code{us54}'s race-to-five win condition, a failed speculative declare is no longer at worst a
lost opportunity but always a two-point swing.

That sign change is what motivated the roster's threshold arithmetic. Writing $p$ for a
speculative plan's success probability and $q$ for the probability the set gets banked correctly
anyway if the seat waits, the declare-now-versus-wait comparison is
\begin{equation*}
\underbrace{p \cdot (+1) + (1-p)\cdot 0 = p > q}_{\text{\code{pagat48} (void)}}
\qquad\text{versus}\qquad
\underbrace{p \cdot (+1) + (1-p)\cdot(-1) = 2p - 1 > q
  \;\Longleftrightarrow\; p > \tfrac{1+q}{2}}_{\text{\code{us54} (gift)}},
\end{equation*}
so a 48-card risk appetite $t$ maps to a \code{us54} declare threshold
\begin{equation}
t_{\mathrm{us54}} \;=\; \frac{1 + t_{\mathrm{pagat48}}}{2}.
\label{eq:map}
\end{equation}
(The speculative plans this arithmetic governs guess only the intra-team assignment of a set
already certainly the team's --- Section~\ref{sec:mechanism} --- so the \code{pagat48} failure
outcome it models is precisely the void, and the derivation stands on the narrowed rule.)
The roster's thresholds were re-derived through this map rather than ported \cite{roster}:
Balanced $0.80 \to 0.90$, Blitz $0.70 \to 0.85$, Punter $0.55 \to 0.775$, Hoarder
$0.95 \to 0.975$, and the Archivist's separate \emph{foreign} bar $0.90 \to 0.95$. The
\code{declareOnlyWhenCertain} styles also ship bars of their own --- Banker $0.95$, Turtle
$0.98$, Scout at the baseline $0.90$ --- values their certainty gate makes moot but that travel
with the vectors. Note the range this induces: \textbf{every shipped style's bar sits at
$0.775$ or above}, the shipped span being $[0.775, 0.98]$. That band is the protagonist of
this paper.

\subsection{Style is not skill}

The laboratory's design rule, stated before any style was written \cite{botlab}: ``every style
shares one identical, full-strength knowledge engine.'' Styles differ only in the policy layer
above \code{buildKnowledge()}; inference depth, memory window, and error rate live on a separate
\emph{skill} vector. The rationale is attribution --- if an aggressive bot loses, one must be able
to say whether aggression is bad or whether the aggressive bot was merely written worse. The same
separation is what makes the instrument of Section~\ref{sec:instrument} meaningful: two style
vectors handed the same view and seed differ only by policy, so any divergence is attributable
to the knobs and nothing else.

\subsection{The \code{StyleParams} surface}

A style is a flat vector of roughly two dozen knobs \cite{styles}: ask-scoring weights
(\code{wHit}, \code{wProgress}, \code{wNarrow}, \code{certaintyBonus}, \code{gambleBonus},
\code{minHitP}); declare policy (\code{declareThreshold}, \code{declareMaxUncertain},
\code{declareOnlyWhenCertain}, \code{declareOnlyOwnHand}); window and clinch behaviour new under
\code{us54} (\code{declareEagerness}, \code{foreignDeclare}, \code{foreignDeclareThreshold},
\code{clinchAggression}, \code{denialWeight}); information policy, tempo targeting, and two
hoarding knobs (\code{hoardBooks}, \code{minHandSize}). The nine roster entries --- Balanced,
Blitz, Punter, Banker, Turtle, Hoarder, Scout, Ghost, Archivist --- are named presets over this
vector, and their aggressive-to-conservative story is carried, nominally, by
\code{declareThreshold}.

\section{The instrument: decision-level divergence}
\label{sec:instrument}

The probe is simple to state and exact in what it certifies. Real \code{us54} games are played
on the laboratory's own seed set. At \emph{every} decision point of every game, both policy
vectors under comparison are handed the identical \code{SeatView} --- the public observation a
human client would receive --- and the identical seed, and their two \code{GameAction}s are
compared structurally. The game itself advances under one designated vector; the other is a
shadow evaluated counterfactually at each point. The statistic is the count (and rate) of
decision points at which the two actions differ.

Three properties of the policy layer make this well-defined \cite{decide}:

\begin{itemize}
  \item \textbf{Purity.} \code{decide(view, policy, seed)} is a pure function: same view, policy,
  and seed imply the same action. The only randomness is a seeded generator over \code{seed}.
  \item \textbf{Shared inference.} Both vectors run the identical full-strength inference engine,
  so the comparison isolates the policy layer by construction.
  \item \textbf{Exactness.} Because the comparison is of actions, not outcomes, a result of zero
  is a \emph{certificate}: the two vectors are the same player on this population. There is no
  standard error to report, and none is needed.
\end{itemize}

The contrast with an outcome-level sweep is the methodological point of this paper. A sweep that
varies a knob and reports win rates over $N$ games can, at best, bound the knob's effect by a
confidence interval around zero. The divergence probe answers a strictly prior question --- did
the knob change any decision at all? --- and answers it exactly, at the cost of replaying
decisions rather than whole extra tournaments. Outcome differences require decision differences;
a knob with zero divergence needs no games spent on it.

\section{The finding}
\label{sec:finding}

\subsection{A roster entry indistinguishable from the control}

The axis was caught because one style failed harder than flat: it vanished. In the pilot payoff
matrix (36 cells $\times$ 200 duplicate-deal pairs), the Hoarder --- the roster's optionality
style, bar at $0.975$, plus the two hoarding knobs --- scored a mean of $0.5484$, identical to
Balanced's $0.5484$ \emph{in every digit}. Its entire row of the payoff matrix and its entire
row of the diagnostics table were byte-identical to Balanced's; the head-to-head cell
\code{balanced-vs-hoarder} measured $0.5000$ with a standard error of $0.0000$; and the
$\alpha$-Rank evaluator had to be given a tie tolerance to stop ranking the two styles on
floating-point noise \cite{styles, alpharank}. A roster entry that cannot be distinguished from
the control is not a style, so the identity was instrumented directly rather than inferred from
the score.

\begin{table}[ht]
\centering
\caption{The divergence instrument on Hoarder versus Balanced, before and after the repair of
Section~\ref{sec:repair}. ``Before,'' the two vectors were the same player: zero divergent
decisions in 400 games.}
\label{tab:beforeafter}
\begin{tabular}{lrrr}
\toprule
Hoarder vs.\ Balanced & decisions & different actions & rate \\
\midrule
before & 275{,}380 & 0 & 0.0000\% \\
after  & 289{,}560 & 9{,}315 & 3.2169\% \\
\bottomrule
\end{tabular}
\end{table}

\subsection{Attribution, one knob at a time}

Zero divergence over 275{,}380 decisions (Table~\ref{tab:beforeafter}) says the whole vector is
inert; it does not say why. Attribution applied each of the Hoarder's seven behavioural deltas
to Balanced \emph{alone} --- one knob at a time, over 76{,}972 decisions on identical views and
seeds --- and every single-knob variant returned zero divergence. All seven knobs were dead
individually, which rules out any story in which two deltas cancel.

\subsection{The threshold sweep and the dead band}

Sweeping \code{declareThreshold} on Balanced explains the zero
(Table~\ref{tab:pilotsweep}). Below the roster's range, the knob is alive: at a threshold of
$0$, 892 of 51{,}420 decisions change. Divergence decays as the threshold rises and reaches
exactly zero at $0.775$ --- and stays zero at $0.8$, $0.9$, and $1$. In this pilot population,
setting \code{declareThreshold} to $1$, or setting \code{declareOnlyWhenCertain}, or setting
\code{declareMaxUncertain} to $0$, all changed nothing --- the same statement three ways: the
speculative-declare routine, \code{evClaim}, never returned a plan that survived its gates.

\begin{table}[ht]
\centering
\caption{Pilot sweep: decisions changed (of 51{,}420) when \code{declareThreshold} is varied on
Balanced, against identical views and seeds. Every shipped bar in the roster ($[0.775, 0.98]$)
sits in the zero region, whose measured extent runs through threshold $1$.}
\label{tab:pilotsweep}
\small
\begin{tabular}{l*{11}{r}}
\toprule
\code{declareThreshold} & 0 & 0.5 & 0.55 & 0.6 & 0.65 & 0.7 & 0.75 & \textbf{0.775} & 0.8 & 0.9 & 1 \\
\midrule
decisions changed & 892 & 842 & 414 & 274 & 140 & 18 & 18 & \textbf{0} & 0 & 0 & 0 \\
\bottomrule
\end{tabular}
\end{table}

The immediate casualty list is wider than one knob. Five parameters ---
\code{declareThreshold}, \code{declareThresholdStalled}, \code{declareEagerness},
\code{clinchAggression}, and \code{denialWeight} --- are consumed only inside \code{evClaim}, so
every one of them is inert for \emph{any} style whose bar sits at or above $0.775$. Blitz
($0.85$), Punter ($0.775$), and the three \code{declareOnlyWhenCertain} styles are all affected;
they remain distinguishable from the control only through their \emph{ask} knobs. The knobs are
wired, swept, and reported; they are not reached.

\subsection{The axis is dead for the whole roster}

After the repair of Section~\ref{sec:repair} and at the full-precision matrix stage, the sweep
was re-run on the shipped roster: \code{declareThreshold} varied on Balanced alone, over 40
\code{us54} games, comparing every decision against the identical view and seed
(Table~\ref{tab:rostersweep}) \cite{styles}.

\begin{table}[ht]
\centering
\caption{Roster-stage sweep, 40 \code{us54} games: divergent decisions when
\code{declareThreshold} is varied on Balanced. Every roster style's bar sits at $0.775$ or
above.}
\label{tab:rostersweep}
\begin{tabular}{lrr}
\toprule
\code{declareThreshold} & divergent decisions & rate \\
\midrule
0.50 & 350 & 1.2968\% \\
0.65 & 48 & 0.1779\% \\
0.70 & 29 & 0.1075\% \\
0.75 & 19 & 0.0704\% \\
\textbf{0.775} & \textbf{0} & \textbf{0.0000\%} \\
0.85 / 0.90 / 0.95 / 0.975 & 0 & 0.0000\% \\
\bottomrule
\end{tabular}
\end{table}

Every roster style sits at $0.775$ or above, so \textbf{the axis is dead for all nine of them}.
The styles themselves are not: measured with the same instrument, Scout differs from Balanced on
$2.89\%$ of decisions, Ghost $2.15\%$, Archivist $1.66\%$, Banker $1.19\%$, Turtle $0.83\%$,
Punter $0.47\%$, Hoarder $0.47\%$, and Blitz $0.39\%$. The roster is genuinely diverse --- just
not along the axis its labels advertise. Punter did not top the payoff matrix by declaring at
$0.775$ rather than $0.90$; it won through \code{gambleBonus} and \code{declareMaxUncertain}.

\section{The mechanism: bimodal confidence}
\label{sec:mechanism}

The pilot audit's first-order reading was ``the speculative path is unreachable: \code{evClaim}
never returns.'' The roster-stage measurement sharpened that, and the sharper statement is the
interesting one: \textbf{the path fires constantly.} Over the same 40 games as
Table~\ref{tab:rostersweep}, Balanced made 171 speculative declares to 138 certain ones. What is
unreachable is not the branch but the \emph{band}: the inference engine's confidence estimates
are bimodal --- a plan sits either well above ${\sim}0.975$ or well below ${\sim}0.775$, almost
never between --- so any threshold inside $[0.775, 0.975]$ --- and, because the upper mode sits
at or arbitrarily near $1$, any measured threshold up through $1$, covering the shipped
$[0.775, 0.98]$ span --- selects the identical set of declares.

Why bimodal? The answer is in the declare pipeline's structure \cite{decide}:

\begin{enumerate}
  \item \textbf{How $p$ is built.} For each unresolved set, \code{planClaim} assigns cards with
  known holders to those holders, and assigns each unknown card by \emph{count-consistency}:
  greedily to the candidate teammate with the most remaining unidentified hand slots,
  decrementing a working capacity map so multiple guesses stay jointly consistent with the
  public hand sizes. The plan's $p$ multiplies per-card estimates --- certain-on-team $=1$,
  certain-on-opponent $=0$, uncertain $=$ chosen capacity over total candidate capacity.
  \item \textbf{What \code{evClaim} may even consider.} A speculative plan qualifies only if its
  guessed cards number between $1$ and \code{declareMaxUncertain} (baseline $1$) \emph{and}
  every guessed card's candidates are all teammates --- that is, the set is already
  \emph{certainly the team's}, with only the intra-team assignment in doubt.
  \item \textbf{Certainty preempts.} In both the own-turn flow and the \code{us54} declare
  window, \code{certainClaim} --- which banks any plan with zero uncertain cards --- runs strictly
  before \code{evClaim}. The public information that establishes ``certainly ours'' (asks, hits,
  count exhaustion) is the same information that pins individual holders, so by the time a set
  qualifies for step 2, its cards are usually located outright and the set exits through
  \code{certainClaim} before \code{evClaim} ever prices it.
\end{enumerate}

What survives to \code{evClaim} is therefore polarised. At one pole, plans whose last card is
forced by capacity arithmetic --- the alternative candidates have no unidentified slots left ---
carry $p$ at or arbitrarily near $1$: nominally speculative, effectively certain. These are the
171 declares that fire, and they clear \emph{every} threshold in the band, which is why the band
is behaviourally uniform from above. At the other pole, plans whose last card genuinely cannot
be pinned have two or more live candidate teammates, and $p$ is a ratio of small integer slot
capacities: splits like $1{:}1$, $2{:}1$, and $3{:}1$ give $0.500$, $0.667$, and $0.750$ --- all
below $0.775$. Estimates \emph{inside} the band require one candidate to hold an outlandish
share of the remaining unidentified slots. The project's own regression suite is the best
witness: to manufacture a mid-band speculative declare for testing, it had to construct a
position with a 21-card teammate hand, yielding $p = 21/22 \approx 0.955$ --- a state real games
essentially never produce \cite{styletest}.

One further consequence deserves its own sentence. The derivation of
equation~\eqref{eq:map} is arithmetically correct --- and, on this engine, currently
\emph{unfalsifiable}: every threshold it can produce in the plausible range yields identical
play, so no experiment run on this roster can confirm or refute the mapping \cite{styles}. A
parameter can be right for reasons the system it governs never samples.

\section{Consequences}
\label{sec:consequences}

\subsection{Repairing the one style built on the dead path}
\label{sec:repair}

The Hoarder's two hoarding knobs, \code{hoardBooks} and \code{minHandSize}, originally gated
\code{evClaim} only --- on the argument that refusing \emph{free} points (a certain declare) to
protect an ask-licence would be a giveaway rather than optionality. Given
Section~\ref{sec:finding}, that argument had already conceded the entire style: the branch it
guarded was never taken. It was also the weaker reading of the style's design intent, since the
card a declare actually spends is nearly always a card of a set the seat is \emph{certain} of.
The gate was therefore relocated to cover every declare the style is free to refuse --- the
speculative one, the certain one, and the wholly-in-hand one --- and never a compulsory
(\code{MUST\_DECLARE}) one \cite{decide, styles}.

The relocation brought the style into existence, measurably
(Tables~\ref{tab:beforeafter} and~\ref{tab:diagnostics}): divergence from Balanced went from
$0.0000\%$ to $3.2169\%$, and every diagnostic moved in the direction the style's thesis
predicts --- including its costs. The head-to-head cell moved from $0.5000 \pm 0.0000$ to
$0.5775 \pm 0.0159$ in Balanced's favour, and the Hoarder's pilot mean score fell from $0.5484$
(tied with the control) to $0.4738$: the style, once expressible, turned out to be priced above
its worth under \code{us54}. That is a finding about the strategy, and a separate one; the
methodological point here is that before the repair there was no strategy to measure.

\begin{table}[ht]
\centering
\caption{Per-style diagnostics pooled over all eight pilot cells: the Hoarder before the gate
relocation is byte-identical to the control; after it, every diagnostic moves. Two Balanced
figures coexist in the source: $0.5484$ is the pooled value in the identity comparison that
flagged the inert Hoarder, while this table's Balanced column pools Balanced's own eight cells
($0.5581$). The load-bearing identity is byte-identical play --- zero divergent decisions in
275{,}380 --- and the table reports both figures rather than resolving what the source does
not.}
\label{tab:diagnostics}
\small
\begin{tabular}{lrrr}
\toprule
diagnostic & Balanced & Hoarder before & Hoarder after \\
\midrule
mean score & 0.5581 & 0.5484 (= Balanced) & 0.4738 \\
foreign-declare rate & 0.018 & 0.018 & 0.247 \\
declare latency (window cycles) & 23.39 & 22.90 & 31.02 \\
race losses / game & 0.044 & 0.046 & 0.091 \\
declares / game & 4.01 & 3.98 & 3.67 \\
forced-declare share & 0.074 & 0.076 & 0.153 \\
concede rate & 0.011 & 0.012 & 0.025 \\
dropout rate & 0.539 & 0.521 & 0.443 \\
ask hit rate & 0.568 & 0.569 & 0.545 \\
\bottomrule
\end{tabular}
\end{table}

\subsection{Relabelling rather than hiding}

A finding like Section~\ref{sec:finding} creates an obligation for every page that publishes the
roster: any description of these nine styles as a spectrum from aggressive to conservative
\emph{declaring} is describing a knob that does not fire. The project chose to relabel its
narrative in place rather than quietly rewrite it. The roster table in the design document now
carries an explicit note that its \code{declareThreshold} figures are 48-card \emph{appetites},
not shipped values --- each is mapped through equation~\eqref{eq:map} before it reaches the
engine --- and that read as shipped values they would contradict the measured inertness result,
which the note cites. The measured-caveats section that records the finding is flagged as
mandatory context ``wherever the ranking is published'' \cite{styles}. The correction is
adjacent to the claim it corrects, which is the only place a correction reliably gets read.

\subsection{A regression gate: no inert vectors}
\label{sec:gate}

An axis that died once can die again --- a refactor of the declare pipeline, a change to the
inference layer's pinning behaviour, or a re-tuned threshold could silently push another knob out
of reach. The laboratory therefore carries the finding forward as a test: for each of eight
probe styles spanning the roster's defining settings, over real bot-game positions under both
rule sets plus two hand-crafted declare-window positions, the suite asserts that the probe style
diverges from the baseline \emph{somewhere} --- at least one position where its action differs
from the control's \cite{styletest}. The test's rationale comment states the doctrine this paper
argues for:

\begin{quote}
``A parameter that is wired but never reachable is worse than an absent one: it makes the
Phase-S1 sweep look like it explored a space it never entered.''
\end{quote}

The two crafted positions are themselves evidence for Section~\ref{sec:mechanism}: they exist
because the Hoarder's gate and the eagerness knob ``only bite on a speculative declare, which
real hard games reach too rarely to rely on.'' Where nature does not supply the state, the
regression net constructs it --- the gate tests reachability of the mechanism, deliberately not
the frequency of the state.

\subsection{The search-side view: \code{inertCandidates}}
\label{sec:inertcand}

The same phenomenon is visible from the opposite direction in the laboratory's exploitability
search, which runs coordinate descent over a ladder of roughly sixty single-knob candidate moves
against each style, accepting or rejecting each move by a sequential probability ratio test
(SPRT; $d_0 = 0$, $d_1 = 0.03$, $\alpha = \beta = 0.05$, 24--400 duplicate pairs per candidate)
and re-measuring accepted results at fixed $N$ on unseen seeds \cite{exploit, botlab, wald}. The
artifact records, per style, the number of candidates whose games were \emph{byte-identical} to
the incumbent's --- paired variance exactly zero --- under the field \code{inertCandidates}
\cite{results}. Table~\ref{tab:inert} extracts the counts from the current full-precision
artifact (matrix v2; its engine commit is recorded \code{1667a1d-dirty}, the suffix reflecting
generation from a modified working tree at that commit).

\begin{table}[ht]
\centering
\caption{Inert candidates in the v2 exploitability search: single-knob moves whose games were
byte-identical to the incumbent's (paired variance $= 0$). Shares are computed from the two
recorded counts.}
\label{tab:inert}
\begin{tabular}{lrrr}
\toprule
target style & candidates tried & inert candidates & share \\
\midrule
Balanced & 56 & 21 & 37.5\% \\
Blitz & 56 & 24 & 42.9\% \\
Punter & 59 & 22 & 37.3\% \\
Banker & 60 & 32 & 53.3\% \\
Turtle & 60 & 29 & 48.3\% \\
Hoarder & 60 & 30 & 50.0\% \\
Scout & 58 & 31 & 53.4\% \\
Ghost & 56 & 22 & 39.3\% \\
Archivist & 56 & 24 & 42.9\% \\
\midrule
total & 521 & 235 & 45.1\% \\
\bottomrule
\end{tabular}
\end{table}

Two design decisions make this statistic honest rather than merely cheap. First, the candidate
ladder is deliberately \emph{not} pruned to the knobs the divergence probe found live: ``inert
at Balanced's settings'' is not ``inert,'' and the ladder visits settings the roster never used
--- \code{declareThreshold} $0.50$, for instance, does move play
(Table~\ref{tab:rostersweep}). Pruning would have made the search circular, able only to
rediscover the axis it was pointed at. Second, the inert candidates cost almost nothing: a
candidate that plays a byte-identical game has exactly zero paired variance and is rejected at
the SPRT's minimum sample size. Decision-level identity is what makes the unpruned, honest
ladder affordable --- and the recorded count is what keeps the search's ``no exploit found''
claims calibrated, since a small exploitability value is only as strong as the number of
candidates that were genuinely, rather than nominally, tried.

\section{Lessons for parameter studies of agents}
\label{sec:lessons}

\begin{enumerate}
  \item \textbf{Report reach, not just effect.} An outcome-level sweep cannot distinguish ``this
  knob has no effect'' from ``this knob was never consulted at a value that mattered.'' The two
  conclusions license very different actions --- the first says the axis is unimportant; the
  second says the experiment never entered the space it claims to have explored. Only
  decision-level instrumentation separates them.
  \item \textbf{Divergence is the honest unit, and it is a certificate.} Zero divergence over
  275{,}380 decisions is not a tight confidence interval around zero; it is exact behavioural
  identity on that population. Conversely, when divergence is nonzero, its rate (here $0.39\%$
  to $2.89\%$ across the roster) states precisely how much behavioural surface an outcome
  experiment is actually powered against.
  \item \textbf{Probe decisions before spending games.} Outcome differences require decision
  differences, so a divergence probe is a cheap necessary-condition filter to run before any
  tournament budget is committed. The exploitability search of Section~\ref{sec:inertcand}
  operationalises the same principle statistically: zero paired variance rejects a candidate at
  the minimum sample size. The laboratory's broader discipline --- SPRT for tuning decisions,
  fixed-$N$ for published estimates, because sequential stopping biases effect sizes ---
  composes cleanly with it \cite{botlab}.
  \item \textbf{Attribute one knob at a time, on identical seeds.} A whole-vector comparison
  detects that something is wrong; only single-delta attribution on paired decisions says which
  knob, and rules out cancellation.
  \item \textbf{Do not name a roster after an unmeasured axis.} The naming here preceded the
  measurement, and the names then had to be walked back in public. Where the narrative axis and
  the measured axis disagree, relabel around the knobs that fire --- and publish the caveat
  wherever the ranking is published.
  \item \textbf{Derived parameters can be unfalsifiable.} Equation~\eqref{eq:map} is correct
  arithmetic whose every plausible output yields identical play. A parameter justified by
  derivation rather than tuning should be checked for whether the system ever samples the region
  where its value matters; if not, say so.
  \item \textbf{Encode the finding as a gate.} A one-off audit decays; a regression test that
  asserts every roster vector diverges from the control somewhere (Section~\ref{sec:gate}) makes
  inertness a build failure instead of a rediscovery.
\end{enumerate}

\section{Limitations}
\label{sec:limitations}

\textbf{The certificate is population-relative.} Zero divergence is exact, but only over the
decision distribution the probe visited: games generated by these policies, on this seed set,
under this rule set. A different opponent population could in principle populate the empty
$[0.775, 0.975]$ confidence band --- though the mechanism of Section~\ref{sec:mechanism} argues the
bimodality follows from the estimator's structure, not from the particular opponents.

\textbf{The result is band-relative, not global.} Below $0.775$ the knob is demonstrably alive
(892 changed decisions at threshold $0$; $1.2968\%$ divergence at $0.50$). ``Inert axis'' here
means inert \emph{across the range the roster spans}, which is the range the roster's narrative
was built on.

\textbf{The roster is untuned.} The design document's tuning protocol --- SPRT re-tuning of every
style against the control --- has not been run; the styles measured are the styles as specified.
A re-tuned roster could in principle move bars below $0.775$, at which point the axis would
come alive and the labels would mean something again.

\textbf{One engine, one estimator.} The bimodality is a property of this count-consistency
confidence estimator under this rule set's information structure, with certainty-preemption in
this pipeline order. The specific breakpoints ($0.775$, $0.975$) will not transfer; the failure
pattern --- certainty preempting speculation until a threshold band depopulates --- plausibly
will, in any deduction-heavy domain where the same evidence that qualifies a plan also resolves
it.

\textbf{Sample sizes differ by stage.} The zero-divergence certificate is strongest for the
pilot audit (275{,}380 and 51{,}420 decisions); the roster-stage sweep of
Table~\ref{tab:rostersweep} covers 40 games. Both agree exactly on the breakpoint.

\section{Conclusion}

A nine-style roster was named along an aggressive-to-conservative declare-threshold axis.
Decision-level instrumentation --- exact comparison of two policies' actions on identical
observations and seeds at every decision point of real games --- showed the axis changes nothing
across the roster's entire range, because the shared inference engine's confidence estimates are
bimodal and every threshold in the band selects the identical declares. The knob was wired,
swept, and reported; it was never reached. The repairs were behavioural (a relocated gate),
editorial (a relabelled narrative, corrected where the claims live), and institutional (a
no-inert-vector regression gate, and a search-side count of byte-identical candidates). For
parameter studies of agents generally: outcomes measure whether a knob \emph{matters}; only
decisions measure whether it \emph{exists}. Measure decisions first.

\begin{thebibliography}{10}
\small

\bibitem{styles}
FishAI project.
\emph{STYLES.md --- the play-style roster under \code{us54}}.
Repository document; \S2 (the \code{StyleParams} surface), \S3.1.1 (the Hoarder inertness
audit), \S6.1 (the inert declare-threshold axis), \S6.4 (matrix v2).

\bibitem{botlab}
FishAI project.
\emph{BOT\_LAB.md --- play-style spectrum: research, experimental design, and metrics}.
Repository document; \S1.3 (style is not skill), \S5.5 (SPRT for tuning, fixed-$N$ for
reporting), \S5.7 (exploitability search).

\bibitem{rules}
FishAI project.
\emph{RULES\_US54.md --- the ``US student'' 54-card variant}.
Repository document; \S1 decision table, row 14 (a declare with any error awards the set to the
opposing team).

\bibitem{decide}
FishAI project.
\emph{lib/engine/bots/decide.ts --- the deterministic bot policy layer}.
Source file: \code{planClaim}, \code{certainClaim}, \code{evClaim}, and the hoard gate
\code{withinHoardLimits}.

\bibitem{roster}
FishAI project.
\emph{lib/engine/bots/roster.ts --- the nine \code{us54} styles}.
Source file; header derivation of $t_{\mathrm{us54}} = (1 + t_{\mathrm{pagat48}})/2$.

\bibitem{styletest}
FishAI project.
\emph{tests/bots/style.test.ts --- the \code{StyleParams} regression net}.
Source file; the no-inert-vector gate and the crafted \code{evWindowView} position.

\bibitem{exploit}
FishAI project.
\emph{lib/lab/analysis/exploit.ts --- best-response search}.
Source file; \code{KNOB\_LADDER} and the \code{inertCandidates} field.

\bibitem{results}
FishAI project.
\emph{src/lab/data/style-results.v2.json --- matrix v2 artifact}.
Data file, engine commit recorded \code{1667a1d-dirty};
\code{exploitability[*].candidatesTried} and \code{.inertCandidates}.

\bibitem{v05}
A.~Hsieh.
\emph{FishAI v0.5: measuring play style without skill in a 54-card Literature variant}.
Companion paper, published alongside this one, 2026: the engine, the roster, and the
full-precision tournament whose reading this paper's finding qualifies.

\bibitem{wald}
A.~Wald.
Sequential tests of statistical hypotheses.
\emph{Annals of Mathematical Statistics}, 16(2):117--186, 1945.

\bibitem{alpharank}
S.~Omidshafiei, C.~Papadimitriou, G.~Piliouras, K.~Tuyls, M.~Rowland, J.-B.~Lespiau,
W.~M.~Czarnecki, M.~Lanctot, J.~Perolat, and R.~Munos.
$\alpha$-Rank: multi-agent evaluation by evolution.
\emph{Scientific Reports}, 9:9937, 2019.

\end{thebibliography}

\end{document}
